% @LITE_DESC: LaTeX Beamer presentation template with multiple slide layouts
% @LITE_SCENE: Conference talks, academic presentations, lecture slides
% @LITE_TAGS: latex, beamer, presentation, slides, talk, conference

\documentclass[aspectratio=169, 11pt]{beamer}

\usetheme{Madrid}
\usecolortheme{default}

\usepackage{graphicx}
\usepackage{booktabs}
\usepackage{amsmath}
\usepackage{hyperref}

% ── Metadata ──────────────────────────────────────────────
\title[Short Title]{Full Presentation Title}
\subtitle{Subtitle Goes Here}
\author[A. Author]{Author Name}
\institute[INST]{Institute / University}
\date{\today}

\begin{document}

% ── Title slide ───────────────────────────────────────────
\begin{frame}
    \titlepage
\end{frame}

% ── Outline ───────────────────────────────────────────────
\begin{frame}{Outline}
    \tableofcontents
\end{frame}

% ── Section 1 ─────────────────────────────────────────────
\section{Introduction}

\begin{frame}{Problem Statement}
    \begin{itemize}
        \item First key point
        \item Second key point
        \item Third key point
    \end{itemize}
\end{frame}

\begin{frame}{Motivation}
    \begin{block}{Why does this matter?}
        Explanation of the motivation and significance.
    \end{block}
    \begin{alertblock}{Key Challenge}
        The main difficulty or open problem.
    \end{alertblock}
\end{frame}

% ── Section 2 ─────────────────────────────────────────────
\section{Method}

\begin{frame}{Our Approach}
    \begin{columns}[T]
        \begin{column}{0.5\textwidth}
            \begin{enumerate}
                \item Step one
                \item Step two
                \item Step three
            \end{enumerate}
        \end{column}
        \begin{column}{0.5\textwidth}
            % \includegraphics[width=\textwidth]{diagram.pdf}
            \centering Diagram placeholder
        \end{column}
    \end{columns}
\end{frame}

\begin{frame}{Mathematical Formulation}
    The objective function:
    \[
        \min_{\theta} \frac{1}{N} \sum_{i=1}^{N} \mathcal{L}(f_\theta(x_i), y_i) + \lambda \|\theta\|^2
    \]
\end{frame}

% ── Section 3 ─────────────────────────────────────────────
\section{Results}

\begin{frame}{Experimental Results}
    \begin{table}
        \centering
        \begin{tabular}{lcc}
            \toprule
            Method & Accuracy & Time (s) \\
            \midrule
            Baseline & 85.2\% & 12.3 \\
            Ours & \textbf{93.1\%} & 8.7 \\
            \bottomrule
        \end{tabular}
        \caption{Comparison on benchmark dataset.}
    \end{table}
\end{frame}

\begin{frame}{Visualization}
    \centering
    % \includegraphics[width=0.7\textwidth]{results_plot.pdf}
    \fbox{\parbox{0.6\textwidth}{\centering\large Plot placeholder}}
\end{frame}

% ── Conclusion ────────────────────────────────────────────
\section{Conclusion}

\begin{frame}{Summary \& Future Work}
    \textbf{Contributions:}
    \begin{itemize}
        \item Contribution 1
        \item Contribution 2
    \end{itemize}
    \vspace{1em}
    \textbf{Future directions:}
    \begin{itemize}
        \item Direction 1
        \item Direction 2
    \end{itemize}
\end{frame}

\begin{frame}
    \centering
    \Huge Thank You! \\ \vspace{1em}
    \normalsize Questions? \\ \vspace{0.5em}
    \texttt{email@example.com}
\end{frame}

\end{document}
